\section{黑龙江的围城与多语边界}

阿尔巴津不是一个由后世国界预先决定归属的点。1685年清军围攻该地、俄方撤离；俄方重返后，1686年清军再次围困。《清圣祖实录》的黑龙江军报显示清廷如何理解防务与行动，俄国西伯利亚衙门文件和驻军报告则记录另一套补给、疾病与领土主张。城址考古可以约束围困发生的物质空间，却不能单独裁定伤亡或谈判责任。守军人数、疾病死亡和停战时点必须保留满文、汉文与俄文材料之间的差异。

1689年七月双方在\eventanchor{EQ-1689-NERCHINSK}{尼布楚签约}。满文、拉丁文、俄文文本的存在和签约本身可靠；耶稣会译员参与的多语谈判并没有消除地名、术语和界线含义的不一致。《清史稿·邦交志》可作后出的检索线索，不能代替条约各本与俄国谈判档。条约把阿尔巴津危机暂时制度化，随后仍须由使节往来、地图制作和地方执行来赋予它具体含义。因此，尼布楚不是一条可直接投射成精确现代线的结论，而是双方在军事僵持、地理知识不对称和翻译条件中达成的文本安排。
